% Hardware, software, and the two computers of the Day 8 lab.
\begin{frame}{Hardware and software}
  \begin{columns}[T]
    \begin{column}{0.46\textwidth}
      \small
      \textbf{Hardware for each group}
      \begin{itemize}
        \item Raspberry Pi 5 (8 GB) with the active cooler and the card of
              your group: \code{pi-NN}.
        \item Power supply, 27 W.
        \item Camera Module 3 with its cable.
        \item Laptop in the Wi-Fi \code{edgeai-lab}.
        \item A bottle and a cup.
      \end{itemize}
    \end{column}
    \begin{column}{0.50\textwidth}
      \small
      \textbf{Software}
      \begin{itemize}
        \item On the board: the environment \code{\textasciitilde/yolo}
              with \code{ultralytics} 8.4.171, NCNN, and LiteRT.
        \item On the laptop: the new environment
              \code{\textasciitilde/day08\_env} with the file
              \code{requirements.txt}: the same version of the
              package. It needs 2.7 GB.
        \item A browser on the laptop, for the live image.
      \end{itemize}
    \end{column}
  \end{columns}
  \vskip 0.4em
  \small
  \renewcommand{\arraystretch}{1.15}
  \begingroup\centering
  \begin{tabular}{@{}L{0.24\textwidth}L{0.70\textwidth}@{}}
    Lab files & \code{Labs/day08/} on the laptop,
                \code{\textasciitilde/edgeai/day08/} on the board \\
    Each new terminal on the board &
                \code{cd \textasciitilde/edgeai/day08} and
                \code{source \textasciitilde/yolo/bin/activate} \\
    The camera cable & Do not connect or remove it when the power is on. \\
  \end{tabular}\par\endgroup
\end{frame}

\begin{frame}{Two computers, one detector}
  \begingroup\centering
  \begin{tikzpicture}[font=\footnotesize, >={Stealth[length=5pt]},
      b/.style={draw=coolgray, rounded corners=3pt, fill=boxgray,
                minimum width=1.75cm, minimum height=1.0cm, align=center},
      lab/.style={font=\scriptsize, text=coolgray, align=center}]
    \node[b] (a) at (0,0)    {Dataset};
    \node[b, draw=treegreen] (b) at (2.15,0) {Train\\(Part B)};
    \node[b, draw=darkorange] (c) at (4.3,0)  {Export\\(Part C)};
    \node[b, draw=darkorange] (d) at (6.45,0) {Live\\detection\\(Part C)};
    \node[b, draw=cardinalred] (e) at (8.6,0)  {Frame-rate\\table\\(Part D)};
    \foreach \a/\b in {a/b, b/c, c/d, d/e} { \draw[->] (\a) -- (\b); }
    \node[lab] at (0,-0.95)    {400 images,\\2 classes};
    \node[lab] at (2.15,-0.95) {YOLO11n,\\25 epochs};
    \node[lab] at (4.3,-0.95)  {six files,\\copy with \texttt{scp}};
    \node[lab] at (6.45,-0.95) {camera,\\web page};
    \node[lab] at (8.6,-0.95)  {image size,\\precision};
    \draw[coolgray, dashed] (-1.0,0.75) rectangle (5.32,-1.45);
    \draw[coolgray, dashed] (5.43,0.75) rectangle (9.6,-1.45);
    \node[lab] at (2.15,1.0) {the laptop};
    \node[lab] at (7.5,1.0)  {the Raspberry Pi};
  \end{tikzpicture}\par\endgroup
  \vskip 0.4em
  \small
  \begin{itemize}
    \item Part A is not in the diagram: it runs two pre-trained models on
          the board. It needs no file of the laptop.
    \item The training needs PyTorch, and the export to LiteRT does not run
          on the board. The two steps stay on the laptop.
    \item The download of the dataset and the training run in the
          background. One student of the group can start sections 0 to 2 of
          the notebook during Part A.
  \end{itemize}
  \keypoint{The board gets no PyTorch file of your model. It gets the
    exported files, and each file has one fixed image size.}
\end{frame}
